\chapter{Conclusion and Strategic Roadmap}
\label{chap:conclusion}

\section{Summary of Scientific Findings}
The \textbf{Quantum Resource Optimizer} project presents an academically grounded, fully reproducible implementation of constrained combinatorial optimization for the clean energy transition, developed for the Qiskit Fall Fest 2026 Track 2 hackathon. By transforming the constrained Markowitz portfolio allocation dilemma into a Quadratic Unconstrained Binary Optimization (QUBO) problem and mapping it to a 2-local Ising Cost Hamiltonian, we have established several key scientific results:

\begin{enumerate}
    \item \textbf{Exact Ground-State Equivalence:} We derived analytical expressions for all Ising exchange couplings $J_{ij}$, local fields $h_i$, and constant energy offset $c_0$. In our software implementation, numerical diagonal extraction confirms that $\langle x | \hat{H}_C | x \rangle \equiv C(x)$ across all $2^N$ computational basis states within numerical precision ($\Delta E < 10^{-12}$).
    
    \item \textbf{Constructive Quantum Interference in the Feasible Subspace:} Empirical benchmarking demonstrates that increasing QAOA circuit depth from $p=1$ to $p=2$ expands the Feasibility Yield from $41.65\%$ to $77.00\%$, an improvement of $+35.35$ percentage points over unguided random sampling ($31.25\%$). This confirms that parameterized alternating unitaries constructively amplify probability amplitudes within the valid budget subspace $\binom{N}{B}$.
    
    \item \textbf{Mitigation of Greedy Heuristic Covariance Blindness:} While the classical greedy heuristic fails to reach the global optimum ($\alpha = 0.9649$) because it ignores off-diagonal covariance cross-terms, QAOA naturally incorporates pairwise entangling gates, reliably locating the globally optimal portfolio ($\alpha = 1.0000$).
    
    \item \textbf{Delineation of Penalty Regimes:} Systematic sensitivity sweeping over $\lambda \in [0.5, 5.0]$ identified $\lambda = 2.5$ as the optimal calibration point. Lower values ($\lambda \le 1.0$) suffer from severe infeasible state leakage, while higher values ($\lambda \ge 5.0$) distort the energy landscape, creating narrow basins that risk trapping classical optimizers.
\end{enumerate}

\section{Roadmap Toward Industrial Quantum Utility}
To advance beyond small-scale statevector demonstration toward utility-scale execution on physical quantum hardware, four architectural innovations must be pursued:

\subsection{1. Constraint-Preserving XY Mixers}
The standard QAOA ansatz utilizes transverse single-qubit $X$ mixers, which flip individual qubits independently and freely transition state vectors between feasible and infeasible Hamming weight sectors. Consequently, soft penalty weights $\lambda$ must be introduced, which distort the objective landscape.

A superior paradigm is the Quantum Alternating Operator Ansatz proposed by Hadfield et al. \cite{hadfield2019quantum}. By replacing the transverse field with an $XY$ mixer:
\begin{equation}
\hat{H}_{XY} = \sum_{(i, j) \in \mathcal{E}} \left( X_i X_j + Y_i Y_j \right)
\end{equation}
the mixer operator commutes with the total particle number operator:
\begin{equation}
\left[ \hat{H}_{XY}, \sum_{i=0}^{N-1} Z_i \right] = 0
\end{equation}
When initialized in a Dicke state $|D_B^N\rangle$ (a uniform superposition strictly over bitstrings with Hamming weight $B$), the state evolution is mathematically confined to the $\binom{N}{B}$ feasible subspace throughout all circuit depths $p$. This eliminates the need for penalty terms entirely, guarantees a 100\% Feasibility Yield by design, and shrinks the active Hilbert space from $2^N$ to $\binom{N}{B}$.

\subsection{2. Quantum Error Mitigation (QEM) on Qiskit Runtime}
Executing deep QAOA circuits on physical QPUs requires active error mitigation. Qiskit Runtime Primitives provide modular resilience levels incorporating:
\begin{itemize}
    \item \textbf{Dynamical Decoupling (DD):} Inserting periodic $X_\pi - X_\pi$ pulse sequences during idle qubit intervals to refocus environmental dephasing noise and suppress $T_2$ decoherence.
    \item \textbf{Zero-Noise Extrapolation (ZNE):} Intentionally scaling noise levels via pulse stretching or unitary folding and extrapolating expectation values back to the zero-noise limit.
    \item \textbf{Pauli Twirling:} Converting coherent gate calibration errors into stochastic depolarizing channels that can be reliably modeled and cancelled.
\end{itemize}

\subsection{3. Commutativity-Aware SWAP Synthesis}
Because all $ZZ$ interaction terms in the diagonal Cost Hamiltonian commute with one another ($[Z_i Z_j, Z_k Z_l] = 0$), the order in which two-qubit entangling gates are executed is completely arbitrary. Quantum compilers can exploit this commutativity to synthesize hardware-aligned linear SWAP networks, ensuring that entangling gates are executed whenever two virtual qubits become adjacent during routing, thereby minimizing circuit depth and CNOT overhead on heavy-hex lattices.

\subsection{4. Hybrid Quantum-Classical Decomposition}
For large-scale portfolios ($N > 100$), monolithic quantum optimization is impractical on NISQ processors. Future architectures should adopt hybrid decomposition frameworks (such as quantum-assisted Benders decomposition or column generation). Classical solvers handle linear constraints, integer bounds, and asset pre-filtering, while the quantum processor solves dense, highly correlated binary sub-problems where classical branch-and-bound exhibits worst-case exponential branching.

\section{Concluding Remarks}
The Quantum Resource Optimizer provides a complete, mathematically verified, and transparent submission for Qiskit Fall Fest 2026 Track 2. By combining rigorous QUBO and Ising formulation with modern Qiskit 2.5 primitives, thorough baseline comparisons, and candid appraisal of physical hardware limits, this work illustrates the discipline required to bridge theoretical quantum algorithms with practical systems engineering.
